\section*{अध्याय \ref{ch:g6:latitude-longitude} --- अक्षांश, देशांतर और समय-क्षेत्र}

\begin{solution}{exo:g6:latitude-longitude:1}
``\omterm{def:g6:latitude-longitude:latitude}{अक्षांश} \omterm{def:g3:continents-and-oceans:equator}{भूमध्य रेखा} से नापा जाता है और $0^\circ$ से
$90^\circ$ तक जाता है (उत्तर या दक्षिण); \omterm{def:g6:latitude-longitude:longitude}{देशांतर} \omterm{def:g6:latitude-longitude:longitude}{प्रमुख याम्योत्तर} से नापा जाता है
और $0^\circ$ से $180^\circ$ तक जाता है (पूर्व या पश्चिम)।''
\end{solution}

\begin{solution}{exo:g6:latitude-longitude:2}
\omterm{def:g3:continents-and-oceans:poles}{उत्तरी ध्रुव}: $90^\circ$\,उ। \omterm{def:g3:continents-and-oceans:poles}{दक्षिणी ध्रुव}: $90^\circ$\,द। \omterm{def:g3:continents-and-oceans:equator}{भूमध्य रेखा} पर का कोई
शहर: $0^\circ$।
\end{solution}

\begin{solution}{exo:g6:latitude-longitude:3}
सिडनी: \omterm{def:g3:continents-and-oceans:equator}{दक्षिणी गोलार्ध}, पूर्व। न्यूयॉर्क: \omterm{def:g3:continents-and-oceans:equator}{उत्तरी गोलार्ध}, पश्चिम।
काहिरा: \omterm{def:g3:continents-and-oceans:equator}{उत्तरी गोलार्ध}, पूर्व। रियो डी जनेरो: \omterm{def:g3:continents-and-oceans:equator}{दक्षिणी गोलार्ध},
पश्चिम।
\end{solution}

\begin{solution}{exo:g6:latitude-longitude:4}
(क) सही। (ख) ग़लत: \omterm{def:g6:latitude-longitude:latitude}{अक्षांश रेखाएँ} \omterm{def:g3:continents-and-oceans:poles}{ध्रुवों} की ओर छोटी होती जाती हैं; सबसे
लंबी \omterm{def:g3:continents-and-oceans:equator}{भूमध्य रेखा} है। (ग) सही। (घ) ग़लत: वह लंदन में ग्रीनविच से
होकर जाता है।
\end{solution}

\begin{solution}{exo:g6:latitude-longitude:5}
काहिरा। न्यूयॉर्क। टोक्यो लगभग $36^\circ$\,उ, $140^\circ$\,पू पर है (\omterm{def:g6:latitude-longitude:latitude}{अक्षांश
रेखा} $30^\circ$\,उ से थोड़ा ऊपर, \omterm{def:g6:latitude-longitude:longitude}{याम्योत्तर}
$120^\circ$\,पू से थोड़ा आगे, $150^\circ$\,पू से पहले)।
\end{solution}

\begin{solution}{exo:g6:latitude-longitude:6}
$45 \div 15 = 3$: उनकी दोपहरों में 3~घंटे का अंतर है। दोपहर पहले
शहर~B में होती है, जो ज़्यादा पूर्व में है।
\end{solution}

\begin{solution}{exo:g6:latitude-longitude:7}
$73 \div 15 \approx 4.9$, पूर्णांकित करके 5: क्षेत्र $+5$ है। जब ग्रीनविच में दोपहर
होती है, तब मुंबई में $12 + 5 = 17$, यानी शाम के 5~बजे होते हैं।
\end{solution}

\begin{solution}{exo:g6:latitude-longitude:8}
$30 \times 111 = 3\,330$~km। \omterm{def:g3:continents-and-oceans:poles}{उत्तरी ध्रुव} $90^\circ$ पर है:
$90 \times 111 = 9\,990$~km, लगभग 10\,000~km।
\end{solution}

\begin{solution}{exo:g6:latitude-longitude:9}
$9 - (-5) = 14$~घंटे का अंतर। $8 + 14 = 22$: टोक्यो में सोमवार रात के 10~बजे
हैं।
\end{solution}

\begin{solution}{exo:g6:latitude-longitude:10}
वह सही है। पृथ्वी पश्चिम से पूर्व की ओर घूमती है, इसलिए पूर्व के शहर पहले सूर्य की ओर
ले जाए जाते हैं: वे और पश्चिम के शहरों से पहले उसे उगते, सबसे ऊँचे बिंदु पर पहुँचते
और डूबते देखते हैं।
\end{solution}

\begin{solution}{exo:g6:latitude-longitude:11}
(क) $42 \times 3 = 126$~सेकंड, 2~मिनट से थोड़ा ज़्यादा। (ख)
$126 \div 240 \approx 0.5$: \omterm{def:g6:latitude-longitude:longitude}{देशांतर} का लगभग आधा अंश। (ग)
$0.5 \times 111 \approx 55$~km ($0.525$ के साथ: लगभग 58~km)। लगभग पचास
किलोमीटर की ग़लती: किसी द्वीप को चूक जाने, या किसी ऐसे \omterm{def:g1:landscapes-around-us:water}{तट} से टकरा जाने के लिए काफ़ी
जिसे दूर समझा गया था।
\end{solution}

\begin{solution}{exo:g6:latitude-longitude:12}
(क) $18 - 9 = 9$: ग्रीनविच में सोमवार सुबह के 9~बजे। (ख) $9 + 10 = 19$:
ग्रीनविच में सोमवार शाम के 7~बजे। (ग) $19 - 8 = 11$: लॉस एंजेलिस में सोमवार सुबह के
11~बजे। जहाज़ सोमवार शाम 6~बजे (टोक्यो के समय) चला और
उसी सोमवार सुबह 11~बजे (लॉस एंजेलिस के समय) उतरता है: पहुँचने पर की घड़ियाँ
चलने के समय की घड़ियों से पहले का समय दिखाती हैं, क्योंकि जहाज़ ने पूर्व की ओर
\omterm{def:g6:latitude-longitude:date-line}{तिथि रेखा} पार की और एक दिन पीछे लौट गया।
\end{solution}

\begin{solution}{pb:g6:latitude-longitude:1}
\textbf{1.} एना (रियो डी जनेरो, $23^\circ$\,द) और मिया (ऑकलैंड,
$37^\circ$\,द)।

\textbf{2.} एना और सैम ($43^\circ$\,प और $118^\circ$\,प)। लीना
\omterm{def:g6:latitude-longitude:longitude}{प्रमुख याम्योत्तर} पर है।

\textbf{3.} लीना: लंदन $52^\circ$\,उ पर है, बाक़ी सबसे ज़्यादा \omterm{def:g3:continents-and-oceans:equator}{भूमध्य रेखा} से
दूर।

\textbf{4.} मिया: ऑकलैंड $175^\circ$\,पू पर है, केवल $180 - 175 = 5$
अंश की दूरी पर, $180^\circ$ \omterm{def:g6:latitude-longitude:longitude}{याम्योत्तर} से।

\textbf{5.} $360^\circ \div 24 = 15^\circ$.

\textbf{6.} लंदन: $0$। रियो डी जनेरो: $43 \div 15 \approx 2.9$, पश्चिम,
इसलिए $-3$। टोक्यो: $140 \div 15 \approx 9.3$, पूर्व, इसलिए $+9$। ऑकलैंड:
$175 \div 15 \approx 11.7$, पूर्व, इसलिए $+12$। लॉस एंजेलिस:
$118 \div 15 \approx 7.9$, पश्चिम, इसलिए $-8$।

\textbf{7.} रियो डी जनेरो: $20 - 3 = 17$, शुक्रवार शाम के 5~बजे। टोक्यो:
$20 + 9 = 29$; $29 - 24 = 5$: शनिवार सुबह के 5~बजे।

\textbf{8.} टोक्यो बहुत पूर्व में है: उसकी घड़ियाँ लंदन की घड़ियों से 9~घंटे आगे हैं।
लंदन में रात 8~बजे तक टोक्यो आधी रात पार कर चुका होता है, इसलिए उसका
शनिवार शुरू हो चुका है।

\textbf{9.} $20 - 8 = 12$: शुक्रवार दोपहर 12~बजे।

\textbf{10.} $12 - (-8) = 20$~घंटे आगे।

\textbf{11.} पूर्व की ओर \omterm{def:g6:latitude-longitude:date-line}{तिथि रेखा} पार करते हुए जहाज़ पिछले
दिन में लौट जाता है।

\textbf{12.} $20 + 12 = 32$; $32 - 24 = 8$: कॉल शुरू होते समय ऑकलैंड में
\textbf{शनिवार सुबह के 8~बजे} हैं। मिया नाश्ते के समय कॉल में जुड़ती है, जबकि
सैम शुक्रवार को दोपहर का खाना खा रहा होता है।
\end{solution}
